\chapter{cna-samples and cna-examples}
\label{ch:samples-and-examples}

This book's ecosystem chapter introduced both of these sibling repositories briefly, as two
different kinds of proof that CNA works. This chapter returns to them from a practical
angle: what each one is actually useful for, if you are learning CNA or evaluating it for a
port.

\section{cna-samples: historical fidelity, and a precise measurement of one gap}

\texttt{cna-samples} is a direct C\texttt{++} port of the official Microsoft XNA Game
Studio~4.0 sample collection --- eighty-six samples in the full upstream set --- licensed
Ms-PL specifically because it is derived Microsoft content, unlike \texttt{cna-examples}
below. Building it follows the same sibling-checkout pattern as CNA itself
(Chapter~\ref{ch:building-cna}): \texttt{../cna} and \texttt{../sharp-runtime}
must sit alongside it on disk.

Its real value to a reader of this book is as a precise, checkable measurement of the
\texttt{.fx} bytecode gap from Chapter~\ref{ch:shader-fx-gap}: 23 of the 86 samples fail to
run today, and every one of those 23 failures traces to the same root cause --- a compiled
custom effect the content pipeline cannot load. If you are trying to judge, before starting a
port, whether an existing XNA game is likely to run on CNA without hitting this gap,
\texttt{cna-samples} is the closest thing to a real, checkable answer available: a sample
that runs is evidence the gap does not block content shaped like it; a sample that fails is
further, specific confirmation of exactly which content shapes still do.

\section{cna-examples: a coherent, in-app tour, and a project still early}

\texttt{cna-examples} answers a different question, and is deliberately original,
MIT-licensed work rather than a Microsoft-derived port --- precisely so it can carry a
license \texttt{cna-samples} cannot. Conceptually modeled on JavaFX's own \texttt{Ensemble}
sample browser, it is a single, cross-platform application (desktop, Web, and mobile) that
lets a user navigate from a top-level area (Input, Audio, Devices, Net, Media, 2D Graphics,
3D Graphics, and so on) down through categories to individual, runnable demonstrations, all
inside one binary rather than as eighty-six separate small programs.

At the time of writing, this project is an early skeleton: the application boots, shows its
home menu, and navigates from an area down to a category, but individual demo screens
themselves are not yet implemented, and only the \texttt{EASYGL} graphics backend is
currently targeted. This book states that plainly rather than describing the project by its
aspiration --- it is worth knowing about specifically because its \emph{intended} final shape
(one coherent, navigable catalog covering every area this book's own Part~V
chapters describe) is a genuinely useful thing to watch for as it matures, not because it is
usable as a comprehensive demo catalog today.

\section{Reading them together}

The two repositories are complementary rather than redundant. \texttt{cna-samples} tells you
whether \emph{historical, real XNA content} runs, at the granularity of eighty-six real,
external samples CNA did not design around. \texttt{cna-examples} is intended to tell you
whether \emph{CNA's own capabilities}, area by area, are demonstrable in one coherent
application --- a question \texttt{cna-samples}, being a fixed port of pre-existing content,
cannot answer on its own, since it can only ever demonstrate what the original eighty-six
samples happened to exercise.
